\section{\texorpdfstring{半范数 \(\mathbf{N}_p\)}{半范数 N sub p}}

设 \(p\) 是一个实数 \(\geqslant 1\) ，又设 \(\mathcal{F}^p(\mathrm{X},M)\) 是这样的有限数值函数 \(f\) （定义在 \(\mathrm{X}\) 上）所成的集：它使 \(M(|f|^p)\) 为有限。显然，若 \(g\) 是属于 \(\mathcal{F}^p(\mathrm{X},M)\) 的一个函数且 \(|f| \leqslant |g|\) ，则 \(f\) 也属于 \(\mathcal{F}^p(\mathrm{X},M)\) ；这一注记连同 Minkowski 不等式表明，\(\mathcal{F}^p(\mathrm{X},M)\) 中两个函数之和也属于此集；再注意到 \(M\) 是正齐次的，便知 \(\mathcal{F}^p(\mathrm{X},M)\) 是空间 \(\mathbf{R}^{\mathrm{X}}\) （即定义在 \(\mathrm{X}\) 上的一切有限数值函数所成的空间）的一个向量子空间。

对每个数 \(p > 0\) 及每个有限数值函数 \(f\) （定义在 \(X\) 上），令

\[
\mathrm{N} _ {p} (f) = \left(M (| f | ^ {p})\right) ^ {1 / p};
\]

则有 \(\mathrm{N}_p(\lambda f) = |\lambda |\mathrm{N}_p(f)\) 对每个标量 \(\lambda\) 成立；此外，若 \(p \geqslant 1\) ，则由 (7)，

\[
\mathrm{N} _ {p} (f + g) \leqslant \mathrm{N} _ {p} (f) + \mathrm{N} _ {p} (g),\tag{8}
\]

这就证明 \(N_{p}\) 是向量空间 \(\mathcal{F}^{p}(X,M)\) 上的一个半范数 (TVS, II, §1)。

命题 4. --- 设 p 与 q 是两个实数 \textgreater{} 0 ，并令 \(1/r = 1/p + 1/q\) 。对定义在 X 上的任意有限数值函数 f, g，

\[
\mathrm{N} _ {r} (f g) \leqslant \mathrm{N} _ {p} (f) \mathrm{N} _ {q} (g)\tag{9}
\]

只要 \(\mathrm{N}_p(f)\) 与 \(\mathrm{N}_q(g)\) 有限。

事实上，关系式 (9) 可以写成

\[
M (| f | ^ {r} | g | ^ {r}) \leqslant \bigl (M (| f | ^ {p}) \bigr) ^ {r / p} \bigl (M (| g | ^ {q}) \bigr) ^ {r / q},
\]

这不过是把 Hölder 不等式 (5) 应用于数 \(\alpha = r / p\) 与 \(\beta = r / q\) 以及函数 \(|f|^p\) 与 \(|g|^q\) 所得的结果。

推论. --- 假设 \(M(1) = 1\) ；那么，对每个定义在 X 上的有限数值函数 f，映射 \(p \mapsto \mathrm{N}_{p}(f)\) 在 {]}0, +∞{[} 上是递增的。

把不等式 (9) 应用于 g = 1 的情形，我们看到，\(\mathrm{N}_{r}(f) \leqslant \mathrm{N}_{p}(f)\) 对一切 q \textgreater{} 0 成立；既然由 \(1/r = 1/p + 1/q\) 所定义的数 r 跑遍使 0 < r < p 成立的数集（这里 q 跑遍 \textgreater{} 0 的数集），推论便得到证明。

命题 5. --- 对每个定义在 X 上的有限数值函数 f，1/p (p \textgreater{} 0) 的值中使 \(\mathrm{N}_{p}(f)\) 为有限者所成的集 I 或为空集，或为一个区间；若 I 不缩为一点，则映射 \(\alpha \mapsto \log \mathrm{N}_{1/\alpha}(f)\) 或在 I 上是凸的，或在 I 的内部恒等于 \(-\infty\) 。

设 r 与 s 是两个不同的数 \textgreater{} 0 ，使得 1/r 与 1/s 属于 I；问题归结为证明：若

\[
{\frac {1}{p}} = {\frac {t}{r}} + {\frac {1 - t}{s}},
\]

其中 \(0 < t < 1\) ，则

\[
\log \mathrm{N} _ {p} (f) \leqslant t \cdot \log \mathrm{N} _ {r} (f) + (1 - t) \log \mathrm{N} _ {s} (f),\tag{10}
\]

或者等价地，

\[
\mathrm{N} _ {p} (f) \leqslant \left(\mathrm{N} _ {r} (f)\right) ^ {t} \left(\mathrm{N} _ {s} (f)\right) ^ {1 - t},\tag{11}
\]

由 \(N_{p}\) 的定义，这一关系式可以写成

\[
M (| f | ^ {p}) \leqslant \left(M (| f | ^ {r})\right) ^ {t p / r} \left(M (| f | ^ {s})\right) ^ {(1 - t) p / s}.\tag{12}
\]

令 \(\alpha = tp / r\) ，我们有 \(1 - \alpha = (1 - t)p / s\) ，这由把 \(p\) 定义为 \(t, r, s\) 的函数的那个关系式得出；从而 \(p = \alpha r + (1 - \alpha)s\) 。于是 Hölder 不等式给出

\[
M \big (| f | ^ {r \alpha} | f | ^ {s (1 - \alpha)} \big) \leqslant \big (M (| f | ^ {r}) \big) ^ {\alpha} \big (M (| f | ^ {s}) \big) ^ {1 - \alpha},
\]

这正好就是不等式 (12)。

